%% sections/intro.tex: the opening (problem, benchmark, environment, measures). OWNER: Renzhi. Only Renzhi edits this file.
%% Figures in figures/shared/. Numbers: doc/final/final.tex and rl/PLAN.md (EE599-Project, 2026-11-26).

\tikzset{
  introBox/.style={draw=usccardinal, fill=usccardinal!6, rounded corners=2pt, thick, align=center,
                   font=\footnotesize, text width=26mm, minimum height=11mm, inner sep=2pt},
  introSim/.style={introBox, draw=uscgray, fill=uscgray!8},
  introArr/.style={-{Stealth[length=2mm]}, thick, uscgray},
}

\begin{frame}{Sizing an amplifier}
  \begin{columns}[c]
    \begin{column}{0.55\textwidth}
      \begin{itemize}
        \item A \alert{spec} (gain, bandwidth, power, load) becomes a circuit in two steps:
          \begin{itemize}
            \item pick a \alert{topology}, the arrangement of devices;
            \item \alert{size} it: device dimensions, bias and supply.
          \end{itemize}
        \item Every sizing is checked by one SPICE simulation.
        \item For demanding specs few pass: 0.5--3\% of the design space in our tasks, with every transistor in its
          required operating region.
      \end{itemize}
    \end{column}
    \begin{column}{0.42\textwidth}
      \centering
      \begin{tikzpicture}
        \node[introBox] (s) at (0,0)      {\textbf{Spec}};
        \node[introBox] (t) at (0,-15mm)  {\textbf{Topology}};
        \node[introBox] (z) at (0,-30mm)  {\textbf{Sizing}};
        \node[introSim] (v) at (0,-45mm)  {\textbf{SPICE}\\ pass / fail};
        \draw[introArr] (s) -- (t); \draw[introArr] (t) -- (z); \draw[introArr] (z) -- (v);
        \draw[introArr, usccardinal] (v.east) -- ++(6mm,0) |- node[right, pos=0.25, font=\scriptsize] {try again}
          (z.east);
      \end{tikzpicture}
    \end{column}
  \end{columns}
\end{frame}

\begin{frame}{One design, or the region}
  Optimizers and RL sizing agents return \alert{one} design per spec. The \alert{feasible region}, every sizing that
  meets the spec, tells a designer more:
  \begin{itemize}
    \item \textbf{how safely} it is met: deep inside the region, not on its edge;
    \item \textbf{when to stop}: evidence that no sizing of this topology meets it;
    \item \textbf{what to change}: which target limits the region;
    \item \textbf{what carries over} when the spec is revised.
  \end{itemize}
  \vspace{2mm}
  \begin{block}{Exhaustive maps do not scale}
    5 knobs at 7 levels: 16,807 simulations for one circuit at one load; 8 knobs: 5.8 million.
    A map has to come from a few hundred simulations.
  \end{block}
\end{frame}

\begin{frame}{Our approach: learn where to simulate next}
  \centering
  \begin{tikzpicture}
    \node[introBox] (b)  at (0,0)      {\textbf{Belief}\\ one GP per check};
    \node[introBox] (c)  at (36mm,0)   {\textbf{Candidates}\\ 32 per knob};
    \node[introBox] (p)  at (72mm,0)   {\textbf{Policy}\\ scores each candidate};
    \node[introSim] (s)  at (72mm,-20mm) {\textbf{Simulate}\\ the chosen design};
    \node[introBox] (m)  at (0,-20mm)  {\textbf{Map}\\ $p(x)$ = chance $x$ meets};
    \draw[introArr] (b) -- (c); \draw[introArr] (c) -- (p); \draw[introArr] (p) -- (s);
    \draw[introArr] (s) -- node[below, font=\scriptsize] {update} (m);
    \draw[introArr] (m) -- (b);
  \end{tikzpicture}

  \vspace{3mm}
  \begin{itemize}\small
    \item A \alert{learned acquisition function}: an RL policy, trained once across many circuits, picks the next
      design so that the map gets accurate fast.
    \item \textbf{From Milestone 1:} finding \emph{one} passing design was close to saturated (greedy needed
      0.3--0.5$\times$ the simulations of the best baseline). So the final project scores \alert{mapping}.
  \end{itemize}
\end{frame}

\begin{frame}{The benchmark}
  \begin{columns}[T]
    \begin{column}{0.52\textwidth}
      \small
      \begin{itemize}
        \item SkyWater 130\,nm, ngspice 41 (pinned); GF\,180 for transfer.
        \item 22 catalogue amplifier circuits, plus 2 drawn by our LLM designer.
        \item Each circuit: a pool of 4,096--16,384 sizings, \alert{every one simulated}, split in two halves:
          \begin{itemize}
            \item explorers choose from the \emph{candidate} half;
            \item maps are scored on the \emph{scoring} half.
          \end{itemize}
        \item Specs at the hard end: medium, thin and impossible.
        \item A \alert{knob ladder}: one amplifier at 8, 15, 22 and 29 knobs.
      \end{itemize}
    \end{column}
    \begin{column}{0.45\textwidth}
      \centering\small
      \begin{tabular}{@{}p{0.7\linewidth}r@{}}
        \toprule
        Task set \texttt{map1} & Tasks \\
        \midrule
        Training: 12 circuits of 6--8 knobs, 1 of 15 & 97 \\
        Validation: 3 circuits, the 15-knob rung & 27 \\
        Test: 7 circuits, the ladder at 8, 22, 29 & 69 \\
        \midrule
        GF\,180: 3 circuits & 23 \\
        \bottomrule
      \end{tabular}

      \vspace{2mm}
      \footnotesize Replaying stored pools: no comparison costs a new simulation.
    \end{column}
  \end{columns}
\end{frame}

\begin{frame}{One environment for every explorer}
  \begin{itemize}
    \item \textbf{Episode:} one task (circuit, spec, load) and a budget of 50--600 simulations; 2 starting designs per
      knob.
    \item \textbf{Each step:} 32 candidates per knob from the candidate half; the explorer picks one to simulate.
    \item \textbf{Belief:} one Gaussian process per spec target and per device region check; $p(x)$ is the chance
      every check passes.
    \item \textbf{One small network} (at most 20k parameters) scores every candidate from shared features, so one policy
      runs on any circuit, knob count, spec and load.
    \item \textbf{Reward \texttt{margin}:} the drop in the map's prediction error over 1,024 reference designs,
      weighted toward the end of the run.
    \item \textbf{Train on stored results; validate also live:} on the 15-knob circuit every pick is a fresh, simulated
      design.
  \end{itemize}
\end{frame}

\begin{frame}{Measures and the bar}
  \begin{columns}[T]
    \begin{column}{0.52\textwidth}
      \small
      \begin{itemize}
        \item \textbf{One shared map estimator}, refit on each method's designs: the map depends only on
          \emph{which} designs it chose.
        \item \textbf{\prederr{}}, lower is better: the headline and the choice.
        \item \textbf{\passfail{}}, higher is better: the guard (designs within 0.5\,dB of a gain target count
          neither way).
        \item \textbf{Selection:} lowest \prederr{} on validation, if the \passfail{} interval against the bar
          reaches 0.
        \item \textbf{The bar:} the \rulename{}, picking the highest $p(1-p)$.
      \end{itemize}
    \end{column}
    \begin{column}{0.45\textwidth}
      \centering\small
      The bar on validation (19 tasks, 2 seeds)

      \vspace{1mm}
      \begin{tabular}{@{}lccc@{}}
        \toprule
        simulations & 50 & 200 & 600 \\
        \midrule
        \prederr{} & 0.492 & 0.201 & 0.117 \\
        \passfail{} & & & 0.674 \\
        \bottomrule
      \end{tabular}

      \vspace{2mm}
      \footnotesize Greedy (highest $p$) is within 0.01 of the rule; random pick trails it by 0.35--0.37 on
      \passfail{}.
    \end{column}
  \end{columns}
\end{frame}
